%%======================================================================
%%  Chapitre 4 : Réalisation du projet — Rapport 1 (Lot A)
%%======================================================================
\chapter{Réalisation du projet — Rapport 1 (Lot A)}

Ce chapitre présente la \textbf{Réalisation Lot A} du même projet \textit{EduPlatform}. Le périmètre couvre principalement le \textbf{backend}, les \textbf{processus administratifs} et l'\textbf{intégration API}.

%%======================================================================
\section{Périmètre du Lot A}

Le Lot A regroupe les tâches techniques suivantes :
\begin{itemize}[noitemsep]
    \item implémentation de la couche backend (entités, repositories, services, contrôleurs) ;
    \item gestion des espaces \textbf{administrateur} et workflows métiers critiques ;
    \item centralisation des erreurs, format de réponse API et fiabilisation des endpoints ;
    \item intégration des services fichiers et des règles de validation côté serveur.
\end{itemize}

%%======================================================================
\section{Implémentation backend Spring Boot}

\subsection{Architecture technique mise en œuvre}

Le backend a été structuré selon un découpage en couches :
\begin{itemize}[noitemsep]
    \item \textbf{Model} : entités JPA représentant les objets métiers ;
    \item \textbf{Repository} : accès aux données via Spring Data JPA ;
    \item \textbf{Service} : logique métier et orchestration des traitements ;
    \item \textbf{Controller} : exposition des endpoints REST consommés par le frontend.
\end{itemize}

\subsection{Stabilisation API}

Les réponses JSON suivent le modèle \texttt{ApiResponse} pour standardiser les retours succès/erreur. Un \texttt{GlobalExceptionHandler} assure une gestion homogène des erreurs techniques et fonctionnelles.

%%======================================================================
\section{Réalisation des modules administratifs}

Le Lot A couvre l'essentiel des modules orientés administration :
\begin{itemize}[noitemsep]
    \item \textbf{AdminDemandes} et \textbf{AdminDemandesStage} : traitement des demandes ;
    \item \textbf{AdminAbsences} et \textbf{AdminJustifications} : suivi et validation ;
    \item \textbf{AdminSoumissions} et \textbf{AdminNotes} : supervision académique ;
    \item \textbf{AdminReclamations} : gestion du workflow de réclamation ;
    \item \textbf{AdminNotifications} : diffusion et approbation des notifications.
\end{itemize}

%%======================================================================
\section{Intégration fichiers et persistance}

Les traitements d'upload/download ont été intégrés côté backend pour :
\begin{itemize}[noitemsep]
    \item ressources pédagogiques ;
    \item soumissions de devoirs ;
    \item pièces justificatives d'absences ;
    \item fichiers attachés aux notifications.
\end{itemize}

Le service de stockage applique un nommage unique (UUID) et une organisation par catégories pour garantir la traçabilité des documents.

%%======================================================================
\section{Répartition des tâches — Rapport 1}

\begin{table}[H]
\centering
\caption{Division des tâches de réalisation — Lot A}
\label{tab:taches_realisation_lot_a}
\begin{tabularx}{\textwidth}{lX}
\toprule
\textbf{Axe} & \textbf{Tâches prises en charge} \\
\midrule
Backend & Entités JPA, repositories, services, contrôleurs REST, validation des données. \\
Administration & Développement des modules admin : demandes, absences, justifications, réclamations, notifications. \\
API & Uniformisation des réponses via \texttt{ApiResponse} et gestion globale des exceptions. \\
Persistance/Fichiers & Upload/download, structuration du stockage, liaison avec les workflows métiers. \\
Qualité technique & Vérification des endpoints critiques et cohérence des statuts métier. \\
\bottomrule
\end{tabularx}
\end{table}

\newpage
